\subsection{Motivation and our approach}\label{sec.motivation}
It is folklore that the strong local performance of SOMs ultimately stems from the Newton step. Thus, to design an accelerated method with excellent local performance, the second-order oracle should have the ability to detect the local geometry near the optimal solution, providing a signal indicating that the iterates may have
entered the quadratic convergence region, allowing a Newton step to be active.
For that sake, one particular choice could be the trust-region (TR) oracle~\cite{More1983,conn2000trust}. At a given point \( x \in \mathbb{R}^n \), the classic TR oracle solves the following subproblem:
\begin{equation}
    \label{eq.classic tr subproblem}
    \begin{aligned}
        \underset{d \in \mathbb{R}^n}{\min}~ & \nabla f(x)^T d + \tfrac{1}{2} d^T \nabla^2 f(x)  d \\
        \text{s.t.}                          & \quad  \|d\| \le r,
    \end{aligned}
\end{equation}
where \( r \) denotes the trust-region radius. We denote $(d,\lambda) = \tr(x, r)$ as the primal-dual solution to this problem. For an overview of TR methods, we refer readers to two excellent monographs \cite{yuan2000review,yuan2015recent}.

Two features make \eqref{eq.classic tr subproblem} particularly appealing. Near a non-degenerate minimizer, the step reduces to the pure Newton direction once the ball constraint becomes inactive, preserving quadratic convergence, which is the capability absent from other second-order oracles. 
Moreover, \tr{} oracles can be solved in $O(n^3 + n^2\log\log(\epsilon^{-1}))$ arithmetic operations \cite{vavasisProvingPolynomialtimeSphereconstrained1990,yenewcomplexityresult1991}, incurring no significant overhead compared to solving a Newton equation.\footnote{In comparison, the cubic Newton oracle requires an $O\left(\log\left(1/\epsilon\right)\right)$-time 1-dimensional search procedure \cite{nesterov2006cubic,cartisChapter8Subproblem2022}.}
On the practical side, the TR oracle has been thoroughly analyzed \cite{jiangHolderianErrorBounds2022,rojasNewMatrixfreeAlgorithm2001}, so that highly competitive solvers have been developed \cite{rojasNewMatrixfreeAlgorithm2001, adachiSolvingTrustregionSubproblem2017, adachiEigenvaluebasedAlgorithmAnalysis2019} for it and its extensions.
To this end, TR oracles have formed the algorithmic core of general-purpose nonlinear programming solvers such as Knitro~\cite{byrd2006k}, COPT~\cite{geCardinalOptimizerCOPT2022}, and PDFO~\cite{ragonneau2024pdfo}. Similar success can be found in other cross-cutting applications ranging from adversarial training~\cite{yao2019trust} to reinforcement learning~\cite{schulman2015trust}. These advantages serve as the motivation to use \tr{} to address the global-local balance in accelerated SOMs.


However, to our best knowledge,  TR methods have not been accelerated yet as
using TR oracles introduces technical challenges. Classical TR methods cannot match the global complexity of other unaccelerated SOMs~\cite{cartisComplexitySteepestDescent2010,curtis2017trust} because the dual variable $\lambda$ associated with the ball constraint in \eqref{eq.classic tr subproblem} is determined a posteriori and requires careful monitoring~\cite{curtis2018concise}. Achieving global acceleration demands relative stability between this multiplier and the step size, which roughly means that $\lambda$ should not vary too quickly relative to the step size, but enforcing such stability can interfere with transitions to Newton steps near the optimum.

To address this challenge, we design, for the first time, the accelerated trust-region methods that leverage a modified TR oracle, which we call the trust-region oracle for acceleration denoted as \ref{eq.newTR}. In particular, at a point $x \in \mathbb{R}^n$, the \ref{eq.newTR} oracle solves
\begin{equation*}
\begin{aligned}
\min_{d \in \mathbb{R}^n}~ & \nabla f(x)^\top d 
    + \tfrac{1}{2} d^\top (\nabla^2 f(x) + \sigma I) d, \\
\text{s.t.} &~ \|d\| \le r.
\end{aligned}
\tag*{(TR$_+$)}
\end{equation*}
The dual variable associated with the ball constraint is denoted by $\lambda$, which can detect local geometric information of the optimal solution. 
The primal regularization $\sigma$ in \ref{eq.newTR} is used to effectively modulate $\lambda$, and maintains the stability required for global acceleration, enabling our accelerated trust-region method to systematically balance fast local convergence with provable global guarantee.
\subsection{Contribution}
This work provides an answer to the question posed earlier: we demonstrate that the compatibility between the global and local efficiency of the SOMs depends on the degree of global acceleration.
Specifically, based on the \ref{eq.newTR} oracles \cite{jiang2024nonconvexityuniversaltrustregionmethod}, we develop two accelerated TR-type methods for problem~\eqref{eq.main problem}. One can simultaneously achieve both global acceleration and quadratic local convergence, and the other has a faster global convergence rate but loses local efficiency.

The first method, \emph{Accelerated Trust-Region Method with Local Detection} (\Cref{alg.accelerated utr}), implements a local detection mechanism (\Cref{alg.local detection}) that automatically identifies and ``dives into'' the quadratic convergence regions. This design yields an improved global worst-case oracle complexity of $\tilde{O}(\epsilon^{-1/3})$, while achieving a local quadratic rate of convergence. To our best knowledge, this is the first accelerated SOM that provably achieves both $\tilde{O}(\epsilon^{-1/3})$ global complexity and local quadratic convergence. 

The second method, \emph{Accelerated Trust-Region Extragradient Method} (\Cref{alg.ms accelerated utr}), is designed to persuit the limits of global performance, reaching the near-optimal global oracle complexity of $\tilde{O}(\epsilon^{-2/7})$ for convex problems, but loses the excellent local convergence. Thus, our analysis and experiments reveal a phase transition in accelerated trust-region type methods: pushing global efficiency to its limits naturally entails a trade-off with local performance. 
Such a phase transition occurs due to the loss of local detection in \Cref{alg.ms accelerated utr} when pursuing the near-optimal global convergence rate,
where the primal regularizer $\sigma$ is set prior to determining the extrapolation point that is used in the \ref{eq.newTR} oracle~(as in \Cref{alg.accelerated utr}).

Our theoretical findings are corroborated by numerical experiments. On a global scale, all accelerated SOMs, including ours and \cite{nesterov2008accelerating,monteiro2013accelerated}, consistently outperform the non-accelerated trust-region method~\cite{jiang2024nonconvexityuniversaltrustregionmethod} and cubic Newton method~\cite{nesterov2006cubic}. Locally, \Cref{alg.accelerated utr} exhibits quadratic convergence, and \atrms{} cannot retain superlinear convergence near the solution, similar to the empirical conclusions in \cite{carmon2022optimal,jiang2024accelerated,huang2024inexact}. These findings again confirm the necessity of local detection.

\subsection{Related works}
We review related works along two main lines: TR-type methods and accelerated SOMs that aim to achieve faster-than-sublinear efficiency. Despite the strong empirical success of TR-type methods, its theoretical analysis mainly focuses on nonconvex problems~\cite{yeSecondOrderOptimization2005,curtis2017trust,curtis2021trust,zhangHomogeneousSecondorderDescent2025,hamad2022consistently,grapiglia2015convergence,hamad2024simplepracticaladaptivetrustregion}, while  
its global complexity analysis for convex problems remains incomplete.
In fact, in convex optimization, the classical TR has worse convergence rates than other unaccelerated mainstream SOMs. This gap was recently closed in~\cite{jiang2024nonconvexityuniversaltrustregionmethod}. However, whether TR oracles can further benefit from acceleration, as achieved for other SOMs, remains an open question.

As discussed earlier, unaccelerated SOMs often exhibit superior local performance near the optimal solution compared with their accelerated counterparts. A common remedy for this issue is to employ restart strategies in accelerated frameworks~\cite{nesterov_lectures_2018}. However, this introduces a new challenge of deciding when to restart since the optimal restart frequency depends on unknown, problem-specific parameters. In contrast, when stronger global regularity conditions (akin to strong convexity) are imposed on the objective function, the restart schedule can be determined in a more structured form~\cite{ostroukhov2020tensor, lin2025perseus, huang2025approximation} or even becomes unnecessary~\cite{jiang2022generalized, marques2022variants}. Under such assumptions, these methods can accommodate global performance, and thus the intrinsic trade-off between global convergence and local efficiency becomes a secondary concern.

\paragraph{Organization of the paper}
The remaining part of this paper is organized as follows. Section~\ref{sec.priliminary} reviews the background on TR-type oracles and the use of estimate sequence techniques in accelerated second-order optimization. In Section~\ref{sec.first variant}, we present the \Cref{alg.accelerated utr}, which achieves an oracle complexity of $\tilde O(\epsilon^{-1/3})$ and exhibits quadratic convergence in high-accuracy regimes. Section~\ref{sec.second variant} introduces the \Cref{alg.ms accelerated utr}, which explores the phase transition of acceleration and attains a near-optimal oracle complexity of $\tilde O(\epsilon^{-2/7})$. Section~\ref{sec.experiments} reports numerical results on practical tasks, confirming the theoretical advantages of our proposed methods.